\chapter{Equilibrium is not a maximum}\label{ch:equilibrium}

\section{What is being kept in balance?}
The body does not seek the highest possible temperature. The clinic does not become more secure merely by filling one tank without regard to the pressure in its pipes or the water left elsewhere. A living arrangement has several conditions of continued function. Some resources can be too scarce; some can be present in unusable places; some processes can run too quickly or consume the very reserve on which tomorrow depends.

This observation gives us a different kind of question from the one posed by a price or a sale. Instead of asking how large one return can become, we ask which physical conditions must be maintained and how a proposed action changes them. The answer cannot be a single universal target imposed on every place. A patient, a reservoir and an electrical network are not the same object. Their relevant variables and acceptable conditions must be declared for the case at hand.

Physiologists use homeostasis for the maintenance of conditions compatible with life through processes of regulation. The word does not mean motionless perfection. Breathing, circulation and temperature regulation continue while a person is well. In an economy, too, a stable stock might conceal a continuous flow of production, use, repair and replacement. A useful physical account must be able to distinguish a functioning flow from an identical-looking snapshot.

\section{Reference, equilibrium and stability}
To measure departure we first need a reference. Suppose a deliberately simple model has two water stores, A and B, with a total of twenty physical units. A reference of ten units in each place can be written
\begin{equation}
 x^*=(10,10).
\end{equation}
The star means that this is a declared comparison state. It does not assert that ten is the correct real-world stock for a clinic, that the state will be reached, or that it will remain in place. In a real design the reference would need a reason connected to use, safety and measurement.

An equilibrium makes a different claim. If a specified law of change maps a state $x$ to its next state $T(x)$, then an equilibrium $\bar x$ satisfies $T(\bar x)=\bar x$. Under a rule that never acts, every state is an equilibrium, including a state in which store B lacks water. A mathematical fixed point is therefore not automatically a healthy condition.

Stability asks a further question: what happens after a disturbance? A system may be at equilibrium but move farther away when perturbed. It may return toward the reference, or it may return only while energy, materials and functioning equipment remain available. None of these behaviours follows from naming the reference. They require a declared dynamics and, ultimately, a test of how that dynamics behaves.

\begin{intuitionbox}{Three claims, three different questions}
A \emph{reference} is where we compare. An \emph{equilibrium} is stationary under a specified rule. \emph{Homeostasis} concerns continued function and response under change. A single calculation may inform all three, but it cannot make them identical.
\end{intuitionbox}

\section{A physically possible reference}
Conservation places a useful limit on the discussion. The two stores together hold twenty units. Lossless transfers between them preserve that total. The reference $(10,10)$ is compatible with the total, whereas $(12,12)$ requires twenty-four units and cannot be reached by internal transfers alone. No better incentive or smarter calculation can supply those missing four units without an explicitly represented input.

Compatibility is not enough. A route may be absent or blocked. A source may lack enough available stock for a proposed transfer. The action may use electricity or lose water between the two measurements. We must therefore separate the question of where we would like a system to be from the question of what its physical network can actually do.

Now compare $(18,2)$ with $(14,6)$. Both contain twenty units. The latter is closer to $(10,10)$ in an ordinary geometric sense. But what does closer mean when two resources have different units, or when a four-unit shortage at B matters differently from a four-unit surplus at A? We need a declared measure of deviation before we can assign a number to improvement. Chapter~\ref{ch:reference} will give each coordinate a reference and a scale; Chapter~\ref{ch:potential} will show why a Gaussian-shaped construction supplies a transparent, smooth measure.

\section{Why the direction must depend on the state}
The same transfer can help in one state and harm in another. Moving four units from A to B changes $(18,2)$ into $(14,6)$, bringing both stores toward the symmetric reference. Starting instead from $(10,10)$, the identical movement produces $(6,14)$, moving both away. An action name such as water delivery cannot carry a permanent positive sign independently of its starting state.

This is precisely why maximising every local transaction is not a satisfactory physical rule. A payment may reward the same delivery at both starting points; a field concerned with the represented balance must distinguish them. It must also account for the whole finite movement, because an action that initially helps can overshoot if carried too far. A local slope tells us a direction near the start, but the value of a completed action must include what happens as the state changes.

There is a caution here that strengthens rather than weakens the proposal. A mathematical measure of departure is not itself a law that opens valves or allocates water. It can inform incentives and expose consequences, but actions still require physical permission and an institution that decides what to do with the information. The problem is to build a trustworthy account first, then test whether a proposed way of using it improves continuing function.

\section{The question that leads to a new account}
Imagine that a community has two possible transfers. Both are feasible. One reduces an immediate shortage but consumes a scarce reserve; another produces less immediate delivery while leaving the stores better prepared for the next disturbance. A monetary invoice records payment, but it does not itself calculate the represented change in physical balance. An endpoint-only stock description can miss the energy carriers, losses and intermediate states needed to describe how either transfer happened.

We want an account whose sign and amount respond to what the physical action actually changes, under a stated reference and boundary. The account should not pretend that a desirable endpoint creates its own water, that a smooth mathematical minimum guarantees recovery, or that one stock coordinate describes every need. But it can make a crucial question answerable: did this accepted finite action reduce the declared burden of the represented system, and by how much?
